\section{Julia and Rust survival components: preserving the clock and the estimand}
\label{sec:exec24-longitudinal-other}

A collections decision is defined by its clock before it is defined by its software language.
Consider accounts eligible for a payment holiday at the start of a month. The bank may
compare initiation within a grace period against continued ordinary servicing, follow
default and prepayment, and value interest receipts during a fixed horizon. A program
that returns a survival curve has solved only one part of this problem. The eligibility
date, treatment strategy, event definition, observation mechanism and sampling unit
must survive the translation from the business request to the estimator. This section
examines six concrete Julia and Rust alternatives at commits preceding the client
cutoff of 7 October 2026. Their role is conditional reuse or reference, not evidence
that a native causal product has already been qualified.

\subsection{A target and an observation contract}
Let $T^a$ denote event time under strategy $a$, measured from eligibility, and let
$S_a(t)=\Pr(T^a>t)$. Two possible targets are the event-risk contrast
\begin{equation}
 \Delta_F(\tau)=\Pr(T^1\leq\tau)-\Pr(T^0\leq\tau),
 \qquad
 \Delta_M(\tau)=\int_0^\tau\{S_1(t)-S_0(t)\}\,dt.
\end{equation}
The first is a probability difference; the second is a restricted-mean survival-time
difference and has time units. If the event is default, prepayment needs an explicit
competing-event or hypothetical-intervention interpretation. Treating prepayment as
ordinary independent censoring changes the question unless the necessary assumptions
are supplied. These definitions do not by themselves identify either counterfactual
curve. Consistency, treatment and censoring exchangeability under the specified
history, positivity and an appropriate observation law remain substantive conditions.

Write $L_i$ for entry, $C_i$ for right censoring and $\widetilde T_i$ for observed exit.
For a continuous clock with a declared boundary convention, a risk indicator might be
$Y_i(t)=\mathbf1\{L_i<t\leq\widetilde T_i\}$. By contrast, an interval-censored
record $(u_i,v_i]$ says that the event occurred somewhere in that interval; it does not
say that observation began at $u_i$. Three otherwise eligible accounts, one entering
after the first observed event, give a first-event increment $1/2$, while ignoring
entry gives $1/3$. This exact example concerns data meaning, not numerical speed.
The independent reference also constructs two deterministic event times, four and
two months, with horizon five. Their restricted means differ by two months while
their event risks at the horizon are both one, so the risk difference is zero.
Multiplying the time scale by thirty changes the mean contrast to sixty days and
leaves the probability contrast unchanged.~\citep{exec24_survivalother_reference}

The proposed command and MCP request must therefore bind population, eligibility
clock, intervention strategy, history availability, entry, censoring, competing events,
horizon, output units and sampling unit before selecting an engine. A convenient
two-column schema is insufficient evidence that two engines implement the same target.

\subsection{Julia alternatives and their distinct contracts}
The JuliaStats Survival.jl snapshot declares version 0.3.3 and Julia 1.10 compatibility,
with a direct MIT licence. Its exported event type contains a time and a Boolean
event/right-censoring indicator. The event-table constructors, Kaplan--Meier component
and Cox fitting component provide useful building blocks for a narrow right-censored
worker. They do not, through that response type alone, encode entry or a treatment
history. The inspected Cox objective uses fractional removal of tied event risk
scores, an Efron-style convention.~\citep{exec24_survivalother_juliastats}

Even this narrow candidate needs explicit boundary fixtures. In the inspected
zero-time helper, the first nonzero index is $i$ and the deletion interval is
$\texttt{start}:(\texttt{start}+i-1)$. For a conventional one-based vector
with times $(0,1,2)$, the abstract trace removes indices one and two, leaving only
time two. If all times are zero, the first-nonzero search has no result and the
deletion branch is not entered. The already-sorted event-vector constructor also
passes its vector into this mutating helper; its comment that the input is unaffected
does not settle aliasing. These are source-local concerns with inspected constructor
and export paths, not reported outcomes from a Julia execution. The safe proposal
route is to reject unsupported zero-time inputs until the intended rule and a native
regression fixture are agreed; silently adding a small epsilon would change the
clock and conceal the question.~\citep{exec24_survivalother_juliastats,exec24_survivalother_reference}

SurvivalAnalysis.jl supplies another architecture. Its pinned version 0.1.0 declares
Julia compatibility \texttt{\^{}1.2} and MIT licensing. Its response hierarchy contains
right-, left- and interval-censored forms, but the inspected documentation explicitly
says that only right-censoring methods are included. In particular,
\texttt{Surv(start,stop)} represents interval censoring. The parametric fitting
dispatch takes a right-censored response and admits Weibull or exponential baselines.
The Hessian inverse and standard-error calculation have exception paths producing
NaNs. An adapter must report those outcomes as unavailable uncertainty, not replace
them with zero standard errors. PH prediction distinguishes a linear predictor from
a transformed survival distribution, which is useful for a typed output contract.
The stated candidate role is an original Julia parametric worker, subject to baseline,
optimisation and uncertainty checks; its response classes do not justify a universal
censoring claim.~\citep{exec24_survivalother_juliaanalysis}

SurvivalModels.jl, at its pinned 0.1.6 declaration with Julia 1.11 compatibility and
MIT licensing, develops parametric general-hazard alternatives and Cox and
nonparametric components. The inspected CoxDefault implementation groups tied times
before subtracting $d\log R$ from the log likelihood: this is a Breslow-style tied
denominator. It should not be expected to agree numerically with the preceding
Efron-style path for every tied data set. With risk scores $(1,2,3)$ and the first
two subjects failing together, the numerator is two and the total risk score is six.
The Breslow partial-likelihood contribution is $2/6^2=1/18$, whereas fractional
removal gives $2/[6(6-3/2)]=2/27$. A parity fixture must name its tie convention
before classifying a difference as a failure.~\citep{exec24_survivalother_juliamodels,exec24_survivalother_reference}

Its Kaplan--Meier constructor forms event increments at stored event/censor times.
The vector survival prediction takes their cumulative product; scalar prediction
uses only times strictly less than the requested time. For one event at time two
among two accounts, with the other censored at four, the abstract boundary values
are $S(2-)=1$ and $S(2)=1/2$. The inspected scalar and stored-grid routes thus encode
different sides of the jump, despite documentation describing the usual less-than-or-equal
product. The proposed wrapper must explicitly distinguish left limits from survival
after the event and test scalar/vector agreement under the chosen meaning. A package
name or a successful fit cannot resolve this discrepancy. The selected SurvivalBase and SurvivalDistributions delegates are examined
below. Their observation and distribution contracts qualify this candidate
further; they do not resolve the preceding boundary discrepancy.
\citep{exec24_survivalother_juliamodels,exec24_survivalother_reference}

JointSurvivalModels.jl offers a different reuse candidate: constructing distributions
from hazards linked to a longitudinal process. Its version 1.0.0 snapshot declares
Julia 1 compatibility and MIT licensing. The inspected structure combines a baseline
hazard with covariate and linked-process effects. This can help formulate a stipulated
joint model of servicing history and event time, but a model of association does not
identify a payment-holiday intervention. Numerical support defaults to a positive
lower endpoint $10^{-4}$ and upper endpoint $10^4$; hazard integration begins at
the configured lower endpoint. The sampling method accepts a random generator but
the inspected uniform draw calls the unqualified global random function. Integration
tolerance, tail support and generator routing must therefore be inspected and tested
before reproducibility or distributional accuracy is claimed. A conditional modelling
worker is plausible; a blanket port would import these unresolved semantics.
\citep{exec24_survivalother_juliajoint}

\subsection{What the Julia delegates preserve, and what the model must identify}
\label{rel22:survival-delegates}

The survival proposal becomes more precise when the two named delegates are
read alongside their caller. The selected SurvivalBase revision declares
version 0.1.0 in June 2026; SurvivalDistributions declares 0.1.1 in July 2026.
Both carry direct MIT notices. These are pre-cutoff source snapshots, not an
observed dependency resolution or a tested installation of SurvivalModels.
Their narrower roles are valuable: SurvivalBase constructs formula terms,
and SurvivalDistributions supplies distribution families and hazard operations.
Neither supplies the treatment assignment, censoring independence or
transport assumptions needed for an intervention recommendation.
\citep{rel22survbase,rel22survdist,rel22survcaller}

SurvivalBase's \texttt{Surv} term produces two columns: observed time and an
event indicator. Its README explicitly limits the current binding to
right-censored observations. The selected term has no entry-time field or
pair of interval-censoring endpoints. The parent general-hazard likelihood
then uses the event contribution minus cumulative hazard from origin. This
composition is coherent for its stated observation type, but passing two
convenient columns cannot encode a delayed-entry experiment. Under independent
entry and censoring, conditional on the modeled covariates, an observation
entering at time $l$ and exiting at time $t$ instead contributes
\begin{equation}
 \ell=\delta\log h(t)-H(t)+H(l),\qquad t>l,
 \label{rel22:entry-likelihood}
\end{equation}
where $h$ is the hazard, $H$ its cumulative integral and $\delta$ indicates an
observed event. The last term conditions on survival to entry. It is not
recovered by relabeling entry as an interval-censoring endpoint or by retaining
only the exit time. \citep{rel22survbase,rel22survcaller,rel22survwitness}

An exact small construction shows the consequence for risk costing. Two
accounts enter at month four; one defaults at month five and the other is
right censored at month six. In a stipulated exponential model with constant
rate $\lambda$, the conditional exposure is three account-months and the
maximum-likelihood rate is $1/3$. Treating both accounts as observed from
origin gives eleven account-months and a rate of $1/11$. For a currently
surviving account, an assumed USD~100 loss on default gives next-month
expected losses of USD~28.35 and USD~8.69, respectively, using
$100(1-e^{-\lambda})$. These are calculations under the constructed
observation law, not observed bank losses or fitted Julia outputs. They show
why a right-censored worker must reject unsupported entry information or
delegate to a likelihood that retains it; an engine's convenient schema
cannot decide the scientific clock. \citep{rel22survwitness}

Even correctly encoded observations need an identifiable parameterization.
The parent general-hazard implementation composes the baseline as
\begin{equation}
 h(t\mid x,z)=h_0(te^{z\alpha})e^{x\beta},\qquad
 H(t\mid x,z)=H_0(te^{z\alpha})e^{x\beta-z\alpha}.
 \label{rel22:gh-composition}
\end{equation}
The second expression follows by changing variables in the hazard integral.
The caller's time and hazard multipliers implement this distinction; simply
multiplying cumulative hazard by the hazard-level factor would omit the
time-change correction. Yet retaining both coefficient blocks in an output
does not prove that the data distinguish them. Take a Weibull baseline with
shape $k>0$, scale $s>0$ and one covariate shared by the two blocks, so $z=x$.
Then
\begin{equation}
 H(t\mid x)=(t/s)^k
       \exp\!\left(x\{(k-1)\alpha+\beta\}\right).
 \label{rel22:weibull-ridge}
\end{equation}
Only the displayed linear combination enters the conditional law. Replacing
the coefficient pair by
$(\alpha+c,\beta-(k-1)c)$ leaves every hazard, survival probability and
right-censored likelihood contribution unchanged. For shape two, the pairs
$(0,\log 3)$ and $(\log 2,\log(3/2))$ are an exact example. One labels the
covariate as a pure hazard-level change; the other assigns part to the time
scale, although they predict the same event distribution.
\citep{rel22survcaller,rel22survwitness}

The parent removes coefficient blocks that are inactive by method definition
before inverting an observed Hessian. That avoids one source of singularity,
but its general-hazard branch retains both blocks and does not remove this
baseline-dependent ridge. With an exponential baseline, an accelerated-hazard
time coefficient disappears altogether because the baseline hazard is
constant. A finite optimizer result, inverse-Hessian request or table headed
\emph{identified parameters} cannot establish separate identification. The
defensible alternatives are to restrict the scientific model, report the
identified combination and its uncertainty, or justify a different baseline
and covariate design with diagnostics for weak identification. No distribution
fit alone turns either coefficient into a causal effect of payment holidays.
These statements concern exact parameter redundancy; no optimizer failure
frequency is claimed. \citep{rel22survcaller,rel22survwitness}

Parameter domains also have to survive the call boundary. The selected
LogLogistic family represents a logistic distribution for log event time.
Writing its location as $\mu\in\mathbb R$ and its positive scale as $\sigma$,
its median is $e^\mu$. The parent's baseline-fitting path exponentiates every
baseline parameter, however, including this location. For a baseline-only fit
without regression covariates, a median of half a month requires
$\mu=\log(1/2)<0$ and is outside that transformed fitting domain. Expressing
the same law in days gives a median of fifteen and location $\log 15>0$.
Changing units should not change whether a physical distribution belongs to
the admitted model. A distribution-specific transformation must distinguish
unrestricted locations from positive scales and shapes. Regression may
absorb some scale changes in other designs, so this example is deliberately
limited to the baseline-only fit. It neither asserts that every baseline is
affected nor predicts the optimizer's returned approximation.
\citep{rel22survdist,rel22survcaller,rel22survwitness}

\subsection{Tail behavior and the meaning of a survival output}
\label{rel22:survival-tail}

The piecewise-constant hazard delegate exposes a second consequential
boundary. It declares that the last supplied hazard continues after the final
finite time interval. Its inspected quantile loop nevertheless returns that
interval's endpoint when the requested cumulative hazard exceeds the
accumulated finite-grid hazard. The cumulative-hazard tail refers to
\texttt{L} although the argument is \texttt{D}; the expectation method takes
\texttt{L} but accesses \texttt{D}. Those names are not bound by the inspected
module. The abstract hazard fallback likewise calls numerical helpers that
the selected module does not import or define. These are specific source
obstacles to admitting these paths, not evidence that every distribution in
the package fails. Fixing a variable name alone would not establish the
intended tail law or qualify the numerical implementation.
\citep{rel22survdist}

For a one-month grid interval with hazard $1/5$ per month continuing forever,
the stipulated law is simply exponential. Its ninetieth percentile is
\begin{equation}
 q_{0.9}=5\log 10\simeq11.513\text{ months},
 \label{rel22:tail-quantile}
\end{equation}
whereas the finite-loop trace returns one month. Survival at that endpoint
is still about 0.819, so the difference cannot be dismissed as a negligible
tail. At a stipulated continuous receipt rate of USD~100 per month until
termination, with no discounting, competing event or cost, expected total
receipts are USD~500. Receipts restricted to the first month are only
USD~90.63. Both are legitimate targets when labeled; silently imposing the
grid endpoint would substitute one for the other. A bank requiring a finite
valuation horizon should specify that horizon and integrate to it. A model
claiming an unbounded mean must establish a proper tail and finite expectation.
The example reconstructs the source branch algebraically and does not run the
downloaded Julia implementation. \citep{rel22survdist,rel22survwitness}

Output names are another source of silent substitution. In the parent,
\texttt{predict\_expected} returns cumulative hazard; the survival method
exponentiates its negative. Cumulative hazard is neither a mean event time nor
an event probability. If it equals two at the chosen horizon, the corresponding
event probability is $1-e^{-2}\simeq0.865$. With an assumed USD~100 loss on
that event, expected loss is USD~86.47, not USD~200. The no-time prediction
also evaluates each training subject at its own observed time, whereas a
portfolio decision generally requires a common future horizon. The adapter
must carry target, units, population and horizon with each array rather than
infer them from a method name. \citep{rel22survcaller,rel22survwitness}

These findings support a conditional comparison, with distinct admission
decisions for right-censored parametric fitting, a piecewise hazard
representation and a general-hazard coefficient model. Preserve the useful
formula and distribution interfaces while testing entry/censoring meaning,
parameter domains, common-horizon prediction, likelihood rank, tail inversion
and time-unit changes against independent laws. Include optimizer convergence
and uncertainty failures in the result: the inspected fit stores the minimizer
and objective without retaining a convergence flag in its model fields.
Native qualification and transitive dependency review remain to be done.
The direct MIT notices identify the selected components' declared terms;
they do not establish scientific calibration or clear every dependency.
For the banking decision, choose the least elaborate model whose assumptions
and identified outputs answer the question, and retain uncertainty about the
causal interpretation separately from uncertainty within a fitted survival law.
\citep{rel22survbase,rel22survdist,rel22survcaller}


\subsection{Rust causal labels require method-level inspection}
The Cameron-Lyons survival repository declares version 2.0.0, edition 2024,
Rust 1.94 and MIT licensing at the pinned commit. These are declarations in the
snapshot, not evidence of a tested toolchain or supported release. Its public Rust
module path re-exports target-trial, marginal-structural-model and TMLE-named
functions. Python registration is delegated and is not established merely by a
Rust export. This explicit distinction matters because an integration proposal must
name the entrypoint whose behaviour is being assessed.~\citep{exec24_survivalother_rustcore}

In the inspected target-trial function, a treated clone initiating within the grace
period receives follow-up $\min(\widetilde T-t_A,\tau)$, while the control clone
uses a clock censored at treatment initiation. With eligibility zero, treatment at
month two, event at month ten and grace period three, the untruncated treated duration
is eight rather than ten months from eligibility. With maximum follow-up five, the
treated duration becomes five but its event flag retains the original event status.
The authored trace is $(5,1)$, while an origin-clock observation administratively
censored at five would be $(5,0)$. This does not establish runtime behaviour of
every wrapper; it establishes why a native admission fixture must bind origin,
grace-period strategy, censoring time and event status jointly.
\citep{exec24_survivalother_rustcore,exec24_survivalother_reference}

The inspected censoring-weight routine ignores its censoring-covariate arguments,
assigns protocol-censored clones a constant based on the trimming option and can
then normalise the weights to mean one. It is not a fitted conditional censoring
model in this path. If two hypothetical conditional observation probabilities are
$1/2$ and $1/4$, their inverse weights are two and four and their mean-normalised
weights are $2/3$ and $4/3$. Giving both records ten and normalising yields one
and one; normalisation cannot create the missing conditioning information.
Likewise, a longitudinal weight's numerator model has a meaning distinct from mean
normalisation. Two treatment probabilities $1/4$ and $1/2$ give an unstabilised
weight eight; numerator probabilities $1/2$ and $1/2$ give a stabilised weight two,
whereas mean normalisation of a single weight gives one.
\citep{exec24_survivalother_rustcore,exec24_survivalother_reference}

The output's name and unit require care. The inspected
\texttt{risk\_difference} field contains a difference between restricted means,
so its unit is time. The main calculation weights clones and uses one hundred
grid points; its subject-resampling bootstrap recreates clones without the inspected
censoring-weight call and uses fifty grid points. Fewer than two bootstrap results
trigger a fixed pair of 0.5 standard errors. The sequential-trial helper skips
small eligible sets and retains successful results from a conditional call, so a
delivery wrapper must preserve attempted, skipped, failed and returned trial counts.
These findings motivate rejection of direct causal admission for these entrypoints
at this snapshot, not rejection of every component in the repository.
\citep{exec24_survivalother_rustcore}

The inspected survival-TMLE function assigns outcome one to observations surviving
beyond a horizon, zero to prior observed events and one half to observations censored
before that horizon. A half is an assumption about an unresolved outcome. With one
known survivor, one known event and two unresolved records among four subjects, the
finite-cohort survival fraction can range from $1/4$ to $3/4$; imputing one half for
each unresolved record chooses $1/2$. It does not identify that value. The inspected
function fits a propensity and an outcome regression but contains no explicit
censoring-model adjustment for that assignment. Its restricted-mean integration starts
at the first supplied horizon, while its reported mean-time standard error is an
unweighted combination of pointwise errors without time widths or cross-horizon
covariances. Rescaling months to days must rescale a time-valued standard error;
the abstract expression without time widths does not do so. A method name cannot
substitute for these checks, and no conclusion here extends to uninspected censoring
or TMLE alternatives elsewhere in the repository.
\citep{exec24_survivalother_rustcore,exec24_survivalother_reference}

\subsection{An interval-data candidate and a justified refusal}
The small Rust turnbull snapshot declares version 0.1.0 and edition 2021. No licence
grant was located in the inspected Cargo declaration or complete root tree, so
adoption or redistribution cannot be presumed. Its response type includes interval
and entry information, but the inspected estimator constructs support only from
uncensored exact-event times and an infinity point; handling truncation and variance
is marked unfinished. With only finite intervals $(1,2]$ and $(2,3]$, this abstract
support construction supplies only infinity. Neither interval contains that point,
so both membership rows have zero sum. The normalisation step then has no positive
denominator. This is a reproducible support witness, not a native Rust test result.
The visible convergence loop also lacks an explicit iteration cap.
\citep{exec24_survivalother_rustturnbull,exec24_survivalother_reference}

The defensible proposal disposition is reference-only pending source, rights and
algorithmic repairs. A qualified interval-censoring worker would need an appropriate
support construction, a declared truncation likelihood, finite-iteration and
nonfinite-value handling, and an uncertainty contract. A zero variance placeholder
must not become a claim of certainty. An interval-only fixture is especially valuable
because ordinary exact-event examples can succeed without exercising the capability
named by the package.

\subsection{Adoption choices, public prototypes and full costs}
The preferred architecture preserves selected original engines behind typed,
versioned workers. JuliaStats is a candidate for narrowly specified right-censored
estimation after boundary qualification; SurvivalAnalysis for its declared parametric
response family; SurvivalModels for explicit tie and prediction contracts; and
JointSurvivalModels for a stipulated joint distribution after numerical and random
generator checks. Rust components can be reference or future worker candidates,
but the inspected causal entrypoints require semantic repair before admission.
The interval-data prototype remains reference-only while its licence and support
questions are unresolved. Whole-library ports are not justified by language
preference alone: a port transfers maintenance, rights, numerical and semantic
obligations without supplying identification.

A generated/public prototype can test the merit of this architecture before an
expensive bank experiment. Register exact-event, all-censored, zero-time, delayed-entry,
interval-only, tied-event, competing-event, administrative-censoring and malformed-input
fixtures; match each fixture to a target-specific oracle. Compare scalar and vector
outputs, time-unit rescaling, row permutations, random-seed routing and complete
estimator resampling. Separately evaluate uncertainty under a declared independent
or clustered sampling law. Finite-law examples establish those exact laws; public
data establish only the observable behaviours and assumptions they actually support.
Neither silently supplies counterfactual bank truth.

The CLI/MCP result should carry the source snapshot, environment and rights profile,
original row and subject identities, target and clock, tie/boundary options, diagnostics,
refusal reason, uncertainty method and complete costs. Preserve unsuccessful requests
and human repairs in the denominator. In a stipulated accounting example, twelve
initiated requests costing USD 36 and producing three accepted numeric deliveries
cost USD 3 per initiation but USD 12 per accepted delivery. Correct refusals should
be scored separately from accepted numeric results; deleting them would distort
both reliability and cost. These amounts are authored examples, not vendor prices
or measured performance.~\citep{exec24_survivalother_reference}

The six selected profiles constitute a contrastive review, not exhaustive
other-language coverage. The Julia discovery tail, Rust alternatives, failed
target-trial search and uninspected delegates remain explicit follow-ups. The method
owner must approve the target and observation law; the adapter owner must preserve
units, failures and identities; the rights owner must clear the exact distribution;
and the evaluation owner must preregister public/generated comparisons. This gives
concrete next actions while leaving native, exhaustive and bank acceptance unestablished.


The chapter's common requirement is to preserve the same account history
through every stage of the calculation. Changing the time origin, event
priority, policy law, averaging population or unit of uncertainty can change
the financial decision while leaving a plausible curve on the screen. The
mathematical references and source-specific failures show which distinctions
a proposed worker must expose. The next treatment compares regimes,
heterogeneous effects and continuous doses using those requirements, rather
than assuming that one longitudinal label makes the methods interchangeable.
